\chapter{نتایج و تفسیر}
\label{chp:results}
\section{مقدمه}
در این فصل، نتایج حاصل از اعمال بهبودهای پیشنهادی بر روی دو الگوریتم \lr{ScSR} و \lr{ZSSR} ارائه و تحلیل می‌شوند. هدف اصلی این ارزیابی، بررسی تعادل میان کاهش زمان اجرا و حفظ کیفیت بازسازی تصویر است. برای سنجش کیفیت بازسازی از دو معیار استاندارد \lr{PSNR}\LTRfootnote{Peak Signal-to-Noise Ratio} و \lr{SSIM}\LTRfootnote{Structural Similarity Index} استفاده شده است. تمامی آزمایش‌ها بر روی چهار مجموعه‌داده استاندارد \lr{Set5}، \lr{Set14}، \lr{BSD100} و \lr{Urban100} در بزرگنمایی ۲ برابر انجام شده‌اند. در ادامه ابتدا به تشریح معیارهای ارزیابی و مجموعه‌داده‌ها پرداخته و سپس نتایج هر یک از الگوریتم‌ها به تفصیل بررسی می‌شود.

\section{معیارهای ارزیابی}

\subsection{\lr{PSNR}}
نسبت بیشینه سیگنال به نویز (\lr{PSNR}) یکی از پرکاربردترین معیارهای کمّی برای سنجش کیفیت تصاویر بازسازی‌شده است. این معیار بر اساس میانگین مربعات خطا (\lr{MSE}) بین تصویر بازسازی‌شده و تصویر مرجع محاسبه می‌شود:
\begin{align}
\text{PSNR} = 10 \cdot \log_{10} \left( \frac{L^2}{\text{MSE}} \right)
\end{align}
که در آن $L$ بیشینه مقدار ممکن پیکسل (برای تصاویر ۸ بیتی برابر ۲۵۵) و \lr{MSE} میانگین مربعات اختلاف پیکسل‌های تصویر بازسازی‌شده و تصویر اصلی است. مقادیر بالاتر \lr{PSNR} نشان‌دهنده کیفیت بهتر بازسازی است.

\subsection{\lr{SSIM}}
شاخص شباهت ساختاری (\lr{SSIM}) معیاری ادراکی‌تر نسبت به \lr{PSNR} است که شباهت ساختاری بین دو تصویر را با در نظر گرفتن سه مؤلفه درخشندگی\LTRfootnote{Luminance}، کنتراست\LTRfootnote{Contrast} و ساختار\LTRfootnote{Structure} می‌سنجد. مقدار \lr{SSIM} در بازه $[0, 1]$ قرار دارد و مقادیر نزدیک‌تر به ۱ نشان‌دهنده شباهت بیشتر به تصویر مرجع است.

\section{مجموعه‌داده‌ها}
\label{sec:datasets}

چهار مجموعه‌داده استاندارد در حوزه افزایش وضوح تصویر برای ارزیابی مورد استفاده قرار گرفته‌اند:

\begin{itemize}
\item \textbf{\lr{Set5}:} شامل ۵ تصویر رنگی استاندارد که یکی از قدیمی‌ترین و پرکاربردترین مجموعه‌داده‌های ارزیابی در این حوزه است.
\item \textbf{\lr{Set14}:} شامل ۱۴ تصویر با تنوع بیشتر در محتوا و بافت نسبت به \lr{Set5}.
\item \textbf{\lr{BSD100}:} شامل ۱۰۰ تصویر از مجموعه‌داده \lr{Berkeley Segmentation Dataset} که تنوع بالایی در صحنه‌ها و بافت‌ها دارد.
\item \textbf{\lr{Urban100}:} شامل ۱۰۰ تصویر از فضاها و معماری شهری که دارای الگوهای هندسی منظم و تکرارشونده است.
\end{itemize}

\section{نتایج بهبود الگوریتم \lr{ScSR}}
\label{sec:scsr_results}

جدول \ref{tbl:scsr_results} نتایج مقایسه عملکرد الگوریتم \lr{ScSR} اصلی و نسخه بهبودیافته پیشنهادی را بر روی چهار مجموعه‌داده نشان می‌دهد.

\begin{table}[!ht]
\centering
\caption{مقایسه عملکرد الگوریتم \lr{ScSR} اصلی و پیشنهادی در بزرگنمایی ۲ برابر}
\label{tbl:scsr_results}
\renewcommand{\arraystretch}{1.5}
\begin{tabular}{|c|c|c|c|c|c|c|}
\hline
\multirow{2}{*}{مجموعه‌داده} & \multicolumn{3}{c|}{روش اصلی} & \multicolumn{3}{c|}{روش پیشنهادی} \\
\cline{2-7}
 & \lr{PSNR} & \lr{SSIM} & زمان (ث) & \lr{PSNR} & \lr{SSIM} & زمان (ث) \\
\hline
\lr{Set5} & \lr{35.0045} & \lr{0.9551} & \lr{1420} & \lr{34.3076} & \lr{0.9436} & \lr{51} \\
\hline
\lr{Set14} & \lr{31.3442} & \lr{0.9172} & \lr{1646} & \lr{30.2580} & \lr{0.9011} & \lr{55} \\
\hline
\lr{BSD100} & \lr{30.1233} & \lr{0.8825} & \lr{1737} & \lr{28.8840} & \lr{0.8799} & \lr{56} \\
\hline
\lr{Urban100} & \lr{28.9751} & \lr{0.8723} & \lr{1861} & \lr{27.6912} & \lr{0.8665} & \lr{57} \\
\hline
\end{tabular}
\end{table}

\subsection{تحلیل زمان اجرا}

بارزترین دستاورد روش پیشنهادی، کاهش چشمگیر زمان اجرا است. همان‌طور که در جدول \ref{tbl:scsr_results} مشاهده می‌شود، زمان اجرا به‌طور میانگین از حدود ۱۶۶۶ ثانیه به حدود ۵۵ ثانیه کاهش یافته است که نشان‌دهنده تسریع تقریباً \textbf{۳۰ برابری} است. این کاهش عمدتاً ناشی از سه عامل است:

\begin{enumerate}
\item \textbf{پیچیدگی پایین‌تر \lr{OMP}:} جایگزینی الگوریتم \lr{Feature-Sign Search} با پیچیدگی $O(K^3)$ توسط الگوریتم \lr{OMP} با پیچیدگی $O(mkK)$ که با مقادیر $m = 36$ و $k = 4$، کاهش محاسباتی قابل توجهی را به‌ارمغان می‌آورد.
\item \textbf{استخراج ویژگی یکباره:} حذف محاسبات تکراری فیلترهای مشتقی برای هر پچ و انجام آن یک‌بار برای کل تصویر.
\item \textbf{استفاده از \lr{float32}:} کاهش حجم داده‌های عددی و افزایش سرعت عملیات ماتریسی.
\end{enumerate}

\subsection{تحلیل کیفیت بازسازی}

از نظر کیفیت بازسازی، روش پیشنهادی افت محدودی را نسبت به روش اصلی نشان می‌دهد. میانگین کاهش \lr{PSNR} در چهار مجموعه‌داده حدود \lr{1.08 dB} و میانگین کاهش \lr{SSIM} حدود \lr{0.009} است. بررسی دقیق‌تر نتایج نکات زیر را آشکار می‌کند:

\begin{itemize}
\item در مجموعه‌داده \lr{BSD100}، کاهش \lr{SSIM} تنها \lr{0.0026} است که نشان‌دهنده حفظ بسیار خوب شباهت ساختاری در تصاویر با بافت‌های متنوع است.
\item بیشترین افت \lr{PSNR} مربوط به مجموعه‌داده \lr{Urban100} با \lr{1.28 dB} است. تصاویر این مجموعه‌داده حاوی لبه‌های تیز و الگوهای هندسی منظمی هستند که بازنمایی دقیق‌تری نیاز دارند و محدود کردن تعداد اتم‌ها در الگوریتم \lr{OMP} بیشترین تأثیر را بر آن‌ها می‌گذارد.
\item کمترین افت \lr{PSNR} مربوط به مجموعه‌داده \lr{Set5} با \lr{0.70 dB} است.
\end{itemize}

\subsection{ارزیابی مصالحه سرعت-کیفیت}

با توجه به تسریع ۳۰ برابری در مقابل افت محدود کیفیت (حدود \lr{1 dB} در \lr{PSNR} و کمتر از \lr{0.01} در \lr{SSIM})، روش پیشنهادی مصالحه‌ای مناسب بین سرعت و کیفیت ارائه می‌دهد. این مصالحه به‌ویژه در کاربردهایی که زمان اجرا عامل حیاتی است، از جمله پردازش بلادرنگ تصویر و افزایش وضوح حجم بالای تصاویر، بسیار ارزشمند است. در حالی که الگوریتم اصلی برای پردازش هر مجموعه‌داده به‌طور میانگین حدود ۲۸ دقیقه زمان نیاز دارد، روش پیشنهادی همین کار را در کمتر از یک دقیقه انجام می‌دهد.


\section{نتایج بهبود الگوریتم \lr{ZSSR}}
\label{sec:zssr_results}

جدول \ref{tbl:zssr_results} نتایج مقایسه عملکرد الگوریتم \lr{ZSSR} اصلی و نسخه بهبودیافته پیشنهادی را نشان می‌دهد.

\begin{table}[!ht]
\centering
\caption{مقایسه عملکرد الگوریتم \lr{ZSSR} اصلی و پیشنهادی در بزرگنمایی ۲ برابر}
\label{tbl:zssr_results}
\renewcommand{\arraystretch}{1.5}
\begin{tabular}{|c|c|c|c|c|c|c|}
\hline
\multirow{2}{*}{مجموعه‌داده} & \multicolumn{3}{c|}{روش اصلی} & \multicolumn{3}{c|}{روش پیشنهادی} \\
\cline{2-7}
 & \lr{PSNR} & \lr{SSIM} & زمان (ث) & \lr{PSNR} & \lr{SSIM} & زمان (ث) \\
\hline
\lr{Set5} & \lr{35.6643} & \lr{0.9572} & \lr{338} & \lr{35.2533} & \lr{0.9560} & \lr{35} \\
\hline
\lr{Set14} & \lr{31.3375} & \lr{0.9141} & \lr{419} & \lr{29.9488} & \lr{0.9120} & \lr{45} \\
\hline
\lr{BSD100} & \lr{29.5112} & \lr{0.8975} & \lr{460} & \lr{29.2157} & \lr{0.8968} & \lr{48} \\
\hline
\lr{Urban100} & \lr{35.6442} & \lr{0.9511} & \lr{357} & \lr{35.2440} & \lr{0.9498} & \lr{36} \\
\hline
\end{tabular}
\end{table}

\subsection{تحلیل زمان اجرا}

مشابه نتایج \lr{ScSR}، بهبود پیشنهادی بر \lr{ZSSR} نیز کاهش چشمگیری در زمان اجرا ایجاد کرده است. زمان اجرا به‌طور میانگین از حدود ۳۹۴ ثانیه به حدود ۴۱ ثانیه کاهش یافته که نشان‌دهنده تسریع تقریباً \textbf{۱۰ برابری} است. این کاهش ناشی از ترکیب چند عامل است:

\begin{enumerate}
\item \textbf{توقف زودهنگام:} سیاست توقف زودهنگام مبتنی بر \lr{MSE} بازسازی، مانع از ادامه تکرارهای بی‌ثمر پس از رسیدن به اشباع یادگیری می‌شود. این سازوکار مهم‌ترین عامل کاهش زمان اجرا است.
\item \textbf{بررسی مکرر‌تر پیشرفت:} کاهش فاصله آزمون از ۵۰ به ۴۰ تکرار باعث تشخیص زودتر اشباع یادگیری و فعال‌سازی سریع‌تر سیاست توقف زودهنگام می‌شود.
\item \textbf{کاهش اندازه برش تصادفی:} کاهش \lr{crop\_size} از $128 \times 128$ به $64 \times 64$ حجم محاسبات هر تکرار را به‌صورت مستقیم کاهش می‌دهد.
\end{enumerate}

\subsection{تحلیل کیفیت بازسازی}

از نظر کیفیت بازسازی، روش پیشنهادی افت بسیار محدودی نسبت به روش اصلی نشان می‌دهد. بررسی جزئیات نتایج نکات قابل توجهی را آشکار می‌کند:

\begin{itemize}
\item \textbf{عملکرد بسیار نزدیک در \lr{SSIM}:} میانگین کاهش \lr{SSIM} در چهار مجموعه‌داده تنها \lr{0.0013} است. به‌ویژه در مجموعه‌داده‌های \lr{BSD100} و \lr{Set5}، کاهش \lr{SSIM} به ترتیب \lr{0.0007} و \lr{0.0012} است که از نظر ادراکی عملاً قابل تشخیص نیست.

\item \textbf{افت محدود \lr{PSNR}:} میانگین کاهش \lr{PSNR} حدود \lr{0.62 dB} است. کمترین افت مربوط به \lr{BSD100} با \lr{0.30 dB} و بیشترین مربوط به \lr{Set14} با \lr{1.39 dB} است. در مجموعه‌داده‌های \lr{Set5} و \lr{Urban100}، افت \lr{PSNR} به ترتیب تنها \lr{0.41 dB} و \lr{0.40 dB} بوده که حاکی از حفظ بسیار خوب کیفیت در این مجموعه‌داده‌ها است.

\item \textbf{عملکرد قابل توجه در \lr{Urban100}:} عملکرد الگوریتم \lr{ZSSR} در مجموعه‌داده \lr{Urban100} به‌طور خاص قابل توجه است. \lr{PSNR} به‌دست‌آمده (\lr{35.64 dB} در روش اصلی و \lr{35.24 dB} در روش پیشنهادی) به‌مراتب بالاتر از سایر مجموعه‌داده‌هاست. علت این عملکرد برتر آن است که مجموعه‌داده \lr{Urban100} از تصاویر فضاها و معماری شهری تشکیل شده که ذاتاً دارای الگوهای تکرارشونده هندسی هستند. این الگوهای تکراری دقیقاً خوراک مناسبی برای الگوریتم \lr{ZSSR} فراهم می‌کنند، زیرا این الگوریتم بر اصل تکرارشوندگی الگوها در مقیاس‌های مختلف داخلی تصویر بنا شده است.
\end{itemize}

\subsection{ارزیابی مصالحه سرعت-کیفیت}

با تسریع تقریباً ۱۰ برابری و افت ناچیز کیفیت (به‌ویژه در \lr{SSIM} با میانگین کاهش تنها \lr{0.0013})، روش پیشنهادی مصالحه بسیار مطلوبی ارائه می‌دهد. در حالی که الگوریتم اصلی \lr{ZSSR} برای هر تصویر به‌طور میانگین حدود ۶ دقیقه زمان نیاز دارد، روش پیشنهادی همین فرآیند را در حدود ۴۰ ثانیه انجام می‌دهد. این کاهش زمان بدون تغییر در معماری شبکه یا فرآیند خودترکیب حاصل شده و صرفاً از طریق مدیریت هوشمندانه‌تر فرآیند آموزش به‌دست آمده است.

\section{مقایسه تطبیقی دو بهبود}

جدول \ref{tbl:comparison} خلاصه‌ای از عملکرد هر دو بهبود پیشنهادی را ارائه می‌دهد.

\begin{table}[!ht]
\centering
\caption{خلاصه مقایسه بهبودهای پیشنهادی بر \lr{ScSR} و \lr{ZSSR}}
\label{tbl:comparison}
\renewcommand{\arraystretch}{1.5}
\begin{tabular}{|c|c|c|c|c|}
\hline
الگوریتم & تسریع & \makecell{میانگین افت \\\lr{PSNR (dB)}} & \makecell{میانگین افت\\\lr{SSIM}} & نوع بهبود \\
\hline
\lr{ScSR} & $\sim$\lr{30x} & \lr{1.08} & \lr{0.009} & تغییر الگوریتم بازنمایی تنک \\
\hline
\lr{ZSSR} & $\sim$\lr{10x} & \lr{0.62} & \lr{0.0013} & مدیریت فرآیند آموزش \\
\hline
\end{tabular}
\end{table}

نتایج نشان می‌دهد که بهبود پیشنهادی بر \lr{ScSR} تسریع بیشتری (۳۰ برابر در مقابل ۱۰ برابر) ارائه می‌دهد اما با افت کیفیت بیشتری همراه است. در مقابل، بهبود پیشنهادی بر \lr{ZSSR} افت کیفیت بسیار کمتری دارد، به‌ویژه در معیار \lr{SSIM} که تنها \lr{0.0013} کاهش یافته است. این تفاوت ناشی از ماهیت متفاوت دو بهبود است: در \lr{ScSR} تغییر اساسی در الگوریتم حل مسئله بهینه‌سازی صورت گرفته (از بهینه‌سازی دقیق به بهینه‌سازی حریصانه)، در حالی که در \lr{ZSSR} تنها فرآیند آموزش بهینه‌تر مدیریت شده و معماری و الگوریتم اصلی بدون تغییر باقی مانده‌اند.

\section*{نتیجه‌گیری}

یافته‌های ارائه‌شده در این فصل نشان می‌دهد که با تغییرات هدفمند در الگوریتم‌های افزایش وضوح تصویر می‌توان به کاهش چشمگیر زمان اجرا دست یافت بدون آنکه کیفیت بازسازی به‌طور قابل توجهی آسیب ببیند.

در مورد الگوریتم \lr{ScSR}، جایگزینی الگوریتم \lr{Feature-Sign Search} با \lr{OMP} به‌عنوان مهم‌ترین عامل تسریع، تسریع ۳۰ برابری را با افت متوسط \lr{1.08 dB} در \lr{PSNR} به ارمغان آورد. این نتیجه نشان می‌دهد که در کاربردهای عملی افزایش وضوح مبتنی بر بازنمایی تنک، استفاده از الگوریتم‌های حریصانه سریع‌تر مانند \lr{OMP} جایگزین مناسبی برای الگوریتم‌های بهینه‌سازی دقیق‌تر اما کندتر است.

در مورد الگوریتم \lr{ZSSR}، ترکیب سیاست توقف زودهنگام با تنظیم ابرپارامترهای آموزش، تسریع ۱۰ برابری را با افت ناچیز \lr{0.0013} در \lr{SSIM} حاصل کرد. این نتیجه بیانگر آن است که بخش قابل توجهی از تکرارهای آموزش در الگوریتم اصلی پس از رسیدن به اشباع یادگیری انجام می‌شدند و عملاً سهم ناچیزی در بهبود کیفیت نهایی داشتند.

هم‌افزایی این بهبودها با اهمیت عملی آن‌ها آشکار می‌شود: هر دو الگوریتم بهبودیافته اکنون قادر به پردازش تصاویر در زمان‌هایی هستند که کاربرد عملی آن‌ها را در سناریوهای واقعی‌تر ممکن می‌سازد. کاهش زمان اجرا از مرتبه دقیقه به مرتبه ثانیه، امکان استفاده از این الگوریتم‌ها در کاربردهایی نظیر پردازش تعداد زیاد تصاویر و سیستم‌های نیمه‌بلادرنگ را فراهم می‌آورد. نتایج این پژوهش، پایه‌ای برای تحقیقات آینده در زمینه بهینه‌سازی الگوریتم‌های افزایش وضوح تصویر فراهم کرده و نشان می‌دهد که با تحلیل دقیق گلوگاه‌های محاسباتی و اعمال تغییرات هدفمند، می‌توان به تسریع‌های قابل توجهی دست یافت.
